\section{A separate consequence for binary complementation}
\label{complement:section}

The binary compiler also transfers an exponential lower bound when
the machine recognizing the complementary language may itself be
nondeterministic. This section uses one external algebraic theorem
from \cite{OpenAIComplement2026}, identified below. The determinization
theorem proved earlier is independent of this imported premise.

\subsection{The external path-diagram bound}

For $m\ge1$, take disjoint sets $D^+,D^-$, each with $m$ labels.
A path diagram is a quadruple of relations
\[
 F\subseteq D^+\times D^+,\quad
 B\subseteq D^-\times D^-,\quad
 L\subseteq D^+\times D^-,\quad
 R\subseteq D^-\times D^+.
\]
The positive sign denotes entry from the left and the negative sign
entry from the right. The relations $F,B$ describe through paths;
$L,R$ describe returns to the side of entry. Product means joining
adjacent boundaries and recording finite paths from an outer entry
to the first outer exit. Expanding and contracting traversals of a
consecutive block proves associativity. The identity has identity
through relations and empty return relations. Write $\mathcal T_m$
for this finite monoid and $z\sqsubseteq w$ for inclusion in all four
relations. Multiplication preserves $\sqsubseteq$ in both arguments.

\begin{center}
\fbox{\begin{minipage}{0.92\linewidth}
\begin{theorem}[Imported order-reversing image bound]
\label{complement:imported}
Let $m\ge1$, let $H$ be a finite set with $h\ge2$ elements, and let
$S_0$ be a submonoid of $\mathcal T_m$. Suppose a surjective
multiplicative map $\phi_0:S_0\to\Rel{H}$ satisfies
\[
 z\sqsubseteq w\quad\Longrightarrow\quad
 \phi_0(w)\subseteq\phi_0(z)
 \qquad(z,w\in S_0).
\]
Then
\begin{equation}\label{complement:external-bound}
 2m\ge2^{\lfloor(h-2)/127\rfloor}.
\end{equation}
Identity preservation of $\phi_0$ is not required.
\end{theorem}
This is the theorem entitled \emph{Order-reversing image bound},
source label \texttt{amp:image}, in Section~2 of
\cite{OpenAIComplement2026}. Its full hypotheses and numerical
conclusion are stated here; its proof by missing recurrent classes
and nested corners is not reproduced in this article.
\end{minipage}}
\end{center}

The parameter $m$ counts labels in each direction. Thus the quantity
in \eqref{complement:external-bound} is the total number $2m$ of
directional labels; it is not the degree of the Brauer matching
representation used in the determinization proof.

\begin{lemma}[Finite acceptance and monotonicity]
\label{complement:representation}
For an $s$-state binary 2NFA $C$ in the endmarker model of this article,
there are a homomorphism $\tau:\bits^*\to\mathcal T_{s+1}$,
fixed endmarker diagrams $\lambda,\rho$, and positive labels $a,b$
such that
\begin{equation}\label{complement:acceptance}
 u\in L(C)\quad\Longleftrightarrow\quad
 (a,b)\in F_{\lambda\tau(u)\rho}.
\end{equation}
Consequently $\tau(u)\sqsubseteq\tau(v)$ implies
\begin{equation}\label{complement:monotonicity}
 xuy\in L(C)\quad\Longrightarrow\quad xvy\in L(C)
 \qquad(x,y\in\bits^*).
\end{equation}
\end{lemma}

\begin{proof}
The finite-computation construction of \cite{OpenAIComplement2026}
has the following exact normalization. Add one new success state.
From every original accepting state, allow a stay move into it; in
the new state sweep right and make one designated artificial exit
past the right endmarker. This exit represents acceptance and is
not an additional move available to the original machine. A finite
computation reaches it precisely when an original finite computation
has reached an accepting state, including an accepting initial state.

Use the resulting $s+1$ states as labels in each direction. For each
cell symbol, including markers, a local diagram edge records any
finite sequence of stay transitions followed by one exit move from
the cell. The entry sign records the side of arrival. Finite
computations split into such visits, and local edge witnesses
concatenate back into finite computations. This gives
\eqref{complement:acceptance} with the marker diagrams fixed and
the ordinary-cell products defining $\tau$. Missing transitions and
infinite nonaccepting computations create no finite successful path.
Finally, inclusion of all four relations is preserved by adjoining
the diagrams for $x,y$ and the markers. The designated accepting
edge therefore persists, proving \eqref{complement:monotonicity}.
\end{proof}

\subsection{Encoded contexts give the required antitone quotient}

Let $H=[h]$ and let $B_h$ be the binary language constructed above,
recognized by a $6h+2$-state 1NFA. For a word $w=R_1\cdots R_t$ over the formal
alphabet $\Rel{H}$, write
\[
 r(w)=R_1\cdots R_t,\qquad
 \enc(w)=\enc(R_1)\cdots\enc(R_t),
\]
and use $r(\eps)=\Id_H$, $\enc(\eps)=\eps$. The compiler proves
\begin{equation}\label{complement:compiler-property}
 \enc(w)\in B_h\quad\Longleftrightarrow\quad r(w)\ne\varnothing.
\end{equation}
The transition tables also define $B_h$ on every other binary word,
so complementation below is taken in the full set $\bits^*$.

\begin{proposition}[Direct antitone quotient]\label{complement:quotient}
If an $s$-state 2NFA $C$ recognizes $\bits^*\setminus B_h$, a
submonoid of $\mathcal T_{s+1}$ admits a surjective unital
homomorphism onto $\Rel{H}$ that reverses inclusion.
\end{proposition}

\begin{proof}
Use $\tau$ from Lemma~\ref{complement:representation} and let
\[
 S_0=\{\tau(\enc(w)):w\in\Rel{H}^{*}\}.
\]
This is the submonoid generated by the diagrams of the individual
encoded relations, including the empty product. Suppose
$\tau(\enc(u))\sqsubseteq\tau(\enc(v))$.
Fix $i,j\in H$ and use the binary contexts
$x=\enc(\Id_{\{i\}})$ and $y=\enc(\Id_{\{j\}})$.
For every relation $A$, the product $\Id_{\{i\}}A\Id_{\{j\}}$
is empty exactly when $(i,j)\notin A$. Hence
\eqref{complement:compiler-property}, followed by complement
acceptance and \eqref{complement:monotonicity}, gives
\[
 (i,j)\notin r(u)\quad\Longrightarrow\quad(i,j)\notin r(v).
\]
Thus $r(v)\subseteq r(u)$. Equality of the two diagrams gives this
inclusion in both directions, so
\[
 \phi_0\bigl(\tau(\enc(w))\bigr)=r(w)
\]
is well-defined. Concatenation proves multiplicativity, the empty
word proves unitality, and every $A\in\Rel{H}$ is the image of
$\tau(\enc(A))$. The implication just proved is the required
order reversal.
\end{proof}

\begin{theorem}[Binary nondeterministic complementation]
\label{complement:binary}
Under the imported premise of Theorem~\ref{complement:imported},
every $s$-state 2NFA recognizing $\bits^*\setminus B_h$ satisfies
\begin{equation}\label{complement:state-bound}
 2(s+1)\ge2^{\lfloor(h-2)/127\rfloor},\qquad
 s\ge\frac12\,2^{\lfloor(h-2)/127\rfloor}-1.
\end{equation}
Consequently, for each $n\ge14$, an $n$-state binary one-way NFA has
a complementary language requiring at least
\begin{equation}\label{complement:n-bound}
 \frac12\,2^{\lfloor(n-14)/762\rfloor}-1
\end{equation}
states in every two-way nondeterministic recognizer.
\end{theorem}

\begin{proof}
Apply Theorem~\ref{complement:imported} to
Proposition~\ref{complement:quotient}, with $m=s+1$.
For the last assertion choose $h=\lfloor(n-2)/6\rfloor$ and add
unreachable states to $B_h$ to obtain exactly $n$ states.
The identity
$\lfloor(h-2)/127\rfloor=\lfloor(n-14)/762\rfloor$
gives \eqref{complement:n-bound}.
\end{proof}

The same homomorphism $\tau$ and its $2(s+1)$ directional labels are
used throughout this argument. In particular, encoded word length
does not introduce a state factor. The external dependency is
specifically Theorem~\ref{complement:imported}; the encoding and
the antitone-quotient argument have been given here. This article
provides no new Lean formalization of that imported theorem or of
this complementation transfer. The compiler's finite exact audits
verify its stated encoding behavior and do not certify the external
path-diagram bound.
